\documentclass{stex}
\libusepackage{naproche}
\libusepackage{copyright}
\libinput{preamble}

\NewDocumentCommand\K{}{\mathbf K}

\title{The Clark-Boolos Paradox}
\author{Marcel Schütz}
\date{2026}

\begin{document}

\begin{smodule}{clark-boolos-paradox.ftl}

\usemodule[libraries/meta]{preliminaries.ftl}

\symdecl*{Clark-Boolos Paradox}
\symdecl*{Comprehension Principle}
\symdef{kprop}{\mathbf K}

\maketitle

\noindent
In \cite{KirchnerBenzmuellerZalta2019} Kirchner, Benzmüller and Zalta present
the \emph{Abstract Object Theory} (AOT), a formal metaphysical
theory that distinguishes between ordinary and abstract
objects.
The authors describe that in a naive version of AOT the \emph{Clark-Boolos
Paradox} arises, a paradox similar to Russell's Paradox in naive set theory.

This formalization reconstructs a fragment of this naive version of AOT in
Naproche and shows how it entails the Clark-Boolos Paradox.

\begin{forthel}
  \begin{signature*}
    Let $x$ be an object.
    $x$ is abstract is a relation.
  \end{signature*}

  Let a property stand for a set.
  Let $x$ exemplifies $F$ stand for $x$ is an element of $F$.

  \begin{signature*}
    Let $F$ be a property and $x$ be an object.
    $x$ encodes $F$ is a relation.
  \end{signature*}
\end{forthel}

The key axiom from which the Clark-Boolos Paradox arises
is the comprehension principle for abstract objects:

\begin{forthel}
  \begin{axiom*}[title=Comprehension Principle,name=Comprehension Principle]
    Let $\Phi$ be a class.
    There exists an abstract object $x$ such that for all properties $F$ $x$ encodes $F$ iff $F \Elem \Phi$.
  \end{axiom*}
\end{forthel}

Similar to Russell's Paradox -- which shows that the
assumption of naive set theory that the extension of
every predicate is a set yields a contradiction -- the
Clark-Boolos Paradox shows that assuming that every
predicate is a property yields a contradiction as
well.\footnote{%
  In this formalization which interprets (unary) predicates
  as classes and properties as sets these two assumptions
  coincide.
}

\begin{forthel}
  \begin{definition*}
    $\kprop$ is the class of all objects $x$ such that $x$ encodes some property $F$ such that $x$ does not exemplify $F$.
  \end{definition*}

  \begin{theorem*}[title=Clark Boolos Paradox,name=Clark-Boolos Paradox]
    Assume that $\kprop$ is a property.
    Then we have a contradiction.
  \end{theorem*}
  \begin{proof}
    Define $\Phi \DefEq \EClass{\kprop}$.
    Consider an abstract object $a$ such that for all properties $F$ $a$ encodes $F$ iff $F \Elem \Phi$.
    Then $a$ encodes $\kprop$ and $a$ does not encode any property that is not equal to $\kprop$.
    Hence $a$ exemplifies $\kprop$ iff $a$ does not exemplify $\kprop$.
    Contradiction.
  \end{proof}
\end{forthel}

\printbibliography
\printlicense[CcByNcSa]{Marcel Schütz (2026)}
\end{smodule}
\end{document}
